\documentclass[12pt]{letter}
\usepackage{geometry}
\geometry{margin=1in}
\usepackage{hyperref}

\begin{document}

\begin{letter}{%
Editor-in-Chief\\
\emph{Computational Economics}}

\opening{Dear Editor,}

I am pleased to submit the manuscript ``Structural Fuzzing for Behavioral Model Validation'' for consideration in \emph{Computational Economics}.

Behavioral models in economics and social science are routinely validated by comparing predictions against observations, but the field lacks systematic methodology for answering deeper structural questions: Which model inputs are necessary?  Which are redundant?  How robust are predictions to parameter perturbation?  What is the minimal perturbation that breaks a prediction?  This paper addresses that gap.

We present \emph{structural fuzzing}, a six-component validation framework adapted from software engineering fuzzing methodology and the Bond Index evaluator robustness metric.  The framework provides:
\begin{enumerate}
    \item \textbf{Dimension enumeration:} Exhaustive search over all input subsets to discover the minimal sufficient structure.
    \item \textbf{Pareto frontier analysis:} Identification of the parsimony--accuracy tradeoff across model complexities.
    \item \textbf{Sensitivity profiling:} Ablation-based importance ranking of every model input.
    \item \textbf{Model Robustness Index (MRI):} A single scalar measuring structural reliability under random perturbation, analogous to the Bond Index for LLM evaluator robustness.
    \item \textbf{Adversarial threshold search:} Discovery of minimal parameter changes that flip predictions, revealing fragile structural dependencies.
    \item \textbf{Compositional testing:} Greedy dimension addition revealing optimal ordering and diminishing returns.
\end{enumerate}

The framework is demonstrated on a geometric economics model that predicts human behavior across game theory and prospect theory.  The structural fuzzing campaign discovers that only 3 of 9 available dimensions are necessary (with money contributing zero sensitivity), achieving 2.70\% MAE across 16 targets with an MRI of 1.39.

This work is relevant to \emph{Computational Economics} because it provides a principled, automated methodology for structural validation of computational behavioral models---complementing existing tools (cross-validation, information criteria, Bayesian model comparison) with capabilities they lack: adversarial robustness testing, compositional analysis, and a standardized robustness scalar.

The manuscript has not been published previously and is not under consideration elsewhere.  No human participants were involved.  The software implementation is publicly available under the MIT open-source license.

Generative AI tools (Claude, Anthropic) were used for manuscript preparation and formatting; all intellectual content, methodology design, and experimental analysis are the author's own.

Thank you for considering this submission.

\closing{Sincerely,\\[2em]
Andrew H. Bond\\
Senior Lecturer, Department of Computer Engineering\\
San Jos\'{e} State University\\
\href{mailto:andrew.bond@sjsu.edu}{andrew.bond@sjsu.edu}\\
ORCID: 0009-0003-2599-6158}

\end{letter}
\end{document}
